% Tipo de documento. En este caso es un art�culo, para folios A4, tama�o de la fuente 11pt y con p�gina separada para el t�tulo
\documentclass[a4paper,11pt,titlepage]{article}

% Carga de paquetes necesarios. OrdenesArticle es un paquete personalizado
\usepackage[spanish]{babel} 
\usepackage{eurosym} %para poder poner el s�mbolo del euro con \euro
\RequirePackage[T1]{fontenc}
\RequirePackage[ansinew]{inputenx} 
\usepackage[spanish,cap,cont,title,fancy]{OrdenesArticle}
\usepackage{lmodern}
\usepackage{graphicx}
\usepackage{hyperref}
\usepackage[usenames,dvipsnames]{color}
\usepackage{color}
\usepackage{ifthen}
\usepackage{longtable}
\usepackage{enumerate}
\hypersetup{bookmarksopen,bookmarksopenlevel=4,linktocpage,colorlinks,urlcolor=blue,citecolor=blue,
						linkcolor=blue,filecolor=blue,pdfnewwindow,
						pdftitle={JXTA - Aplicaciones Peer to Peer basadas en Java},
						pdfauthor={Juan Andrada Romero, Jose Domingo L�pez L�pez},
						pdfsubject={Sistemas para la Colaboraci�n}}


% Macro para definir una lista personalizada 
\newenvironment{milista}%
{\begin{list}{\textbullet}%
{\settowidth{\labelwidth}{\textbullet} \setlength{\leftmargin}{\dimexpr\labelsep+\labelwidth+5pt}
\setlength{\itemsep}{\dimexpr 0.5ex plus 0.25ex minus 0.25ex}
\setlength{\parsep}{\itemsep}
\setlength{\partopsep}{\itemsep}
\addtolength{\topsep}{-7.5pt}
}}%
{\end{list}}

% Macro para insertar una imagen
%       Uso: \imagen{nombreFichero}{Factor escala}{Caption (leyenda)}{Label (identificador para referenciarla)}
% -------------------------------------------------------------------------------------------------------------
\def\imagen#1#2#3#4{
 \begin{figure}[h]
 \begin{center}
   \scalebox{#2}{\includegraphics{#1}}
 \caption {#3}
 \label{#4}
 \end{center}
 \end{figure}
}

\begin{document}

% En las p�ginas de portada e �ndices, no hay encabezado ni pie de p�gina
\pagestyle{empty} 

% Se incluye la portada
\begin{titlepage}
	\begin{center}
  	{\LARGE UNIVERSIDAD DE CASTILLA-LA MANCHA} \\
  	\bigskip
  	{\Large ESCUELA SUPERIOR DE INFORM�TICA} \\
  	\vspace{20mm}
  	\includegraphics[scale=0.45, keepaspectratio]{./images/esi_bw.png} \\
  	\vspace{20mm}
  	{\Huge \textbf{Almacenamiento y Recuperaci�n de la Informaci�n}}\\
  	\vspace{10mm}
  	{\LARGE \textsc{\textbf{M�dulo de importaci�n e indexaci�n de documentos}}}\\
  	\vspace{30mm}
  	\today
  	\vspace{40mm}
  	\begin{flushleft}
  		{\large Juan Andrada Romero}\\
  		\vspace{1mm}
  		{\large Jose Domingo L�pez L�pez}\\		
  	\end{flushleft}
	\end{center}
\end{titlepage}
\clearpage

% Texto del reverso de la portada
%%\mbox{}
%%\vspace{18cm}
%%\begin{small}
% Se ajusta la separaci�n entre p�rrafos
%%\parskip=10pt 

%%\copyright~ 2008/2009 Juan Andrada Romero. Universidad de Castilla La Mancha, Escuela Superior de Inform�tica de Ciudad Real.

%%Se permite la modificaci�n, copia y distribuci�n de este documento, seg�n la licencia de documentaci�n GNU (\url{http://www.gnu.org}).

%%Este documento fue compuesto con \LaTeX{}. Im�genes generadas con Power Point y Gimp.
%%\end{small}

%%\newpage

% En las p�ginas de �ndices y prefacio, se utiliza numeraci�n romana
\pagenumbering{Roman}

% Se crea el �ndice
\tableofcontents
% Se pasa p�gina y se a�ade el �ndice de figuras al �ndice principal
%\clearpage\phantomsection
%\addcontentsline{toc}{section}{\listfigurename}
% Se crea el �ndice de figuras
%\listoffigures
% Si se quiere crear un �ndice de tablas se pondr�a: \listoftables

\newpage

% Se ajusta la separaci�n entre p�rrafos
\parskip=10pt

% Se a�ade el prefacio al �ndice
%\clearpage\phantomsection
%\addcontentsline{toc}{section}{Prefacio}

% Comienza el contenido del documento. Se utilizan n�meros ar�bigos y el encabezado y pie de p�gina personalizado
\pagenumbering{arabic}
\pagestyle{fancy}

\section{Introducci�n}


\clearpage
\section{Conceptos de JXTA}


\subsection{Par (\textit{Peer})}

Un par (\textit{peer)} es cualquier entidad de red que implementa uno o m�s protocolos utilizados en JXTA. As�, un par puede residir en una PDA, un tel�fono m�vil, un PC, etc, trabajando de manera independiente y as�ncrona con respecto a otros pares.

Un par debe publicarse en la red para que pueda ser descubierto por otros pares y puedan establecer una comunicaci�n. Por tanto, cada par va a tener un identificador �nico y una direcci�n de red, conocida esta �ltima como \textit{peer endpoint} y utilizada para establecer una comunicaci�n punto a punto entre dos pares. \\
\indent Normalmente, los pares se configuran para ser descubiertos espont�neamente por otros pares de la red, a partir del concepto de \textit{grupo de pares}, que puede ser, por naturaleza, persistente o transitorio.

Los pares en JXTA pueden clasificarse en tres categor�as:

\begin{milista}
	\item \textbf{\textit{Minimal-Edge Peers}}: pares que implementan s�lo los servicios principales JXTA, utilizando otros pares para que act�en como su proxy a la hora de utilizar otros servicios que no pertenecen al \textit{core}de JXTA. Pares que pertenezcan a esta categor�a pueden ser sensores.
	\item \textbf{\textit{Full-Edge Peers}}: pares que implementan todos los servicios de JXTA y pueden participar en todos sus protocolos. Estos pares conforman la mayor�a de pares en una red JXTA y pueden ser tel�fonos, PC's, servidores, etc.
	\item \textbf{\textit{Super Peer}}: pares que implementan y proveen recursos para soportar el despliegue y operaci�n de una red JXTA. Hay tres funciones principales que debe tener un \textit{Super-Peer}:
	\begin{itemize}
		\item \textbf{Retransmisi�n}: usada para almacenar y reenviar mensajes entre pares que no tienen conectividad directa debido a cortafuegos o redes privadas.
		\item \textbf{Citas}: se ocupa de mantener los anuncios globales sobre donde est�n publicados otros pares y ayuda a los pares que act�an como proxy para buscar esos anuncios. 
		\item \textbf{Proxy}: utilizado por los \textit{minimal-edge peers} para tener acceso a todas las funcionalidades de la red JXTA. El par proxy traduce las peticiones, responde a las consultas y proporciona soporte para la funcionalidad de los \textit{minimal-edge peers}.
	\end{itemize}

\end{milista}


\subsection{Grupo de pares (\textit{Peer Group})}

Un grupo de pares es una colecci�n de pares que comparten un conjunto com�n de servicios o intereses. Los pares se auto-organizan en grupos, donde cada uno de los grupos se identifica por un identificador �nico. Cada grupo establece su propia pol�tica de pertenencia al grupo, que puede variar desde una pol�tica abierta (cualquier par puede asociarse al grupo) hasta la muy segura y protegida (que requiere credenciales para poder ser miembro).

Los pares pueden pertenecer a m�s de un grupo simult�neamente. Por defecto, el primer grupo que se instancia es el \textit{Network Peer Group}, y, posteriormente, todos los pares pueden optar por unirse a otros grupos en cualquier momento.

Los protocolos JXTA describen como los pares pueden publicarse, descubrirse, unirse, y monitorear grupos de pares; sin embargo, dichos protocolos no indican cu�ndo o por qu� se crean los grupos de pares. Por otra parte, unirse a un grupo consiste simplemente en instanciar todos los servicios definidos por ese grupo. \\
\indent Hay varias motivaciones para la creaci�n de grupos de pares:

\begin{milista}
	\item Para crear un entorno seguro, ya que se puede aplicar una pol�tica de seguridad. �sta puede ser tan simple como un intercambio de texto plano con nombre de usuario y contrase�a, o tan sofisticada como criptograf�a de llave publica. As�, los miembros del grupo pueden publicar y acceder a contenido protegido.
	\item Para delimitar el alcance, ya que los grupos que permiten la creaci�n de un dominio local espec�fico. Por ejemplo, los pares pueden agruparse para implementar una red de documentos compartidos o una red de CPUs compartidas. Los grupos de pares sirven para subdividir la red en regiones abstractas proporcionando un mecanismo impl�cito de alcance. De este modo, las fronteras del grupo definen el alcance de la b�squeda cuando se busca algo dentro del grupo.
	\item Para crear un entorno de monitorizaci�n, ya que permiten monitorear un conjunto de pares para un fin especial.
\end{milista}


\subsection{Servicios de red (\textit{Network Services})}

Los pares cooperan y se comunican para publicar, descubrir, e invocar servicios de red. Adem�s, los pares pueden publicar m�ltiples servicios los cu�les, a su vez, son descubiertos a trav�s del \textit{Peer Discovery Protocol}. De este modo, los protocolos JXTA reconocen dos niveles de servicios de red:

\begin{milista}
	\item \textbf{Servicios de par}: un servicio de este tipo es solamente accesible por el par que est� publicando ese servicio. Si ese par fallara, el servicio tambi�n fallar�a. 
	\item \textbf{Servicios de grupo de pares}: un servicio de grupo de pares est� compuesto de una colecci�n de instancias del servicio, cooperando entre s� y funcionando en m�ltiples miembros del grupo. S� falla cualquiera de los pares, el servicio colectivo del grupo no se ver� afectado siempre y cuando el servicio este a�n disponible desde al menos un par miembro. 
\end{milista}


\subsection{Servicios de grupos de pares (\textit{Peer Group Services})}

JXTA define un conjunto b�sico de servicios de grupos de pares, que todos los pares miembros del grupo deben implementar, al igual que los pares que se unan a dicho grupo. As�, para que dos pares puedan interactuar a trav�s de un servicio, deben ser ambos miembros del mismo grupo. Los servicios de grupo que cada uno de los pares debe implementar con el fin de participar en una red JXTA son:
 
\begin{milista}
	\item\textbf{Servicio de Destino (\textit{Endpoint Service})}: el servicio de destino o de \textit{punto final} es utilizado para enviar y recibir mensajes entre pares. Este servicio normalmente implementa el protocolo de encaminamiento al destino (ERP), aunque implementaciones m�nimas proporcionan s�lo una peque�a parte del total de este protocolo.
	\item\textbf{Servicio de Resoluci�n (\textit{Resolver Service})}: el servicio de resoluci�n es utilizado para enviar solicitudes de consulta gen�rica a otros pares. Los pares pueden definir e intercambiar consultas para encontrar cualquier informaci�n que pueda ser necesaria (por ejemplo, encontrar otros pares, determinar el estado de un servicio o la disponibilidad de una tuber�a, etc.).
\end{milista}

Adem�s de los servicios b�sicos anteriores, se pueden definir otros servicios est�ndares en los grupos, aunque �stos no tienen por qu� ser implementados por todos los grupos, sino s�lo por aquellos que los necesiten. Los servicios est�ndares que se encuentran en la mayor�a de grupos de pares son:  
 
\begin{milista}

	\item\textbf{Servicio de Descubrimiento (\textit{Discovery Service})}: el servicio de descubrimiento es usado por los pares miembros para buscar recursos en el grupo, tales como otros pares, otros grupos, servicios y tuber�as. 
	\item\textbf{Servicio de afiliaci�n a un grupo(\textit{Membership Service})}: es utilizado por los pares para identificarse con seguridad dentro de un grupo. La identidad de un par permite a las aplicaciones y los servicios determinar qui�n est� solicitando una operaci�n y decidir si la operaci�n debe permitirse o no. Las aplicaciones pueden realizar su propio control de acceso o pueden utilizar el servicio de acceso.
	\item\textbf{Servicio de Acceso (\textit{Access Service})}: el servicio de acceso es usado para validar las peticiones hechas por un par a otro. El par que recibe la petici�n verifica las credenciales del par que la solicita para determinar si se le permite el acceso. Sin embargo, no todas las acciones dentro del grupo necesitan comprobarse con el servicio de acceso; s�lo aquellas acciones que son limitadas a algunos pares necesitan ser comprobadas.
	\item\textbf{Servicio de Tuber�as (\textit{Pipe Service})}: el servicio de tuber�a es utilizado para crear y manejar conexiones de tuber�a  entre los miembros del grupo.
	\item\textbf{Servicio de Monitoreo (\textit{Monitoring Service})}: el servicio de monitoreo se utiliza para permitir a un par monitorizar a otros miembros del mismo grupo de pares. 

\end{milista}


\subsection{M�dulos}

Los m�dulos de JXTA son una abstracci�n de JXTA de bajo nivel usada para representar cualquier pieza de ''c�digo'' y el API que proporciona ese c�digo para su uso. Los m�dulos son utilizados para implementar en JXTA servicios, transporte de mensajes, etc.
 
Los m�dulos proveen una abstracci�n gen�rica para permitir a un par instanciar una funci�n o servicio como, por ejemplo, cuando se une a un grupo, ya que necesitan implementar los servicios de ese grupo. Los m�dulos permiten la representaci�n y la publicaci�n de funcionalidades independientes de la plataforma, y permite a los pares describir e instanciar cualquier tipo de implementaci�n de una funcionalidad. Por ejemplo, un par puede instanciar una implementaci�n en Java o C de una funcionalidad espec�fica.
 
La posibilidad para describir y publicar (o anunciar) una funcionalidad independiente de la plataforma es esencial para soportar el desarrollo de nuevos servicios de grupo. De hecho, JXTA utiliza los anuncios de m�dulos para auto-describir sus propios servicios.

Un m�dulo incluye los siguientes elementos: 

\begin{milista}
	
	\item \textbf{Clases}: las clases de un m�dulo se utilizan principalmente para anunciar la existencia de una funcionalidad. Cada clase de un m�dulo es identificada por un �nico ID, el \textit{ModuleClassID}.
	
	\item \textbf{Especificaci�n}: la especificaci�n de un m�dulo se utiliza para acceder a un m�dulo, ya que contiene toda la informaci�n necesaria para acceder o invocar dicho m�dulo. Por tanto, la especificaci�n de un m�dulo es un m�todo para proporcionar la funcionalidad que implica ese m�dulo, pudiendo existir diversas especificaciones para un determinada clase. Cada especificaci�n se identifica de manera �nica, a trav�s del \textit{ModuleSpecID}.

	\item \textbf{Implementaci�n}: la implementaci�n de un m�dulo es la implementaci�n de una determinada especificaci�n, pudiendo existir m�ltiples implementaciones para una determinada especificaci�n del m�dulo. Cada implementaci�n del m�dulo contiene el \textit{ModuleSpecID} de la especificaci�n asociada que implementa.

\end{milista}

As�, los servicios JXTA pueden usar los m�dulos para identificar la existencia del servicio (las clases del m�dulo), la especificaci�n del servicio (una especificaci�n del m�dulo), o una implementaci�n del servicio (una implementaci�n del m�dulo). Cada uno de estos componentes tiene un anuncio asociado el cu�l puede ser publicado y descubierto por otro par JXTA.
 
Por ejemplo, el servicio de descubrimiento de JXTA (\textit{Discovery Service}) tiene un �nico \textit{ModuleClassID}, identific�ndolo como un servicio de descubrimiento, que es su funcionalidad abstracta. Adem�s, puede tener m�ltiples especificaciones del servicio, utilizando diferentes estrategias para realizar el descubrimiento, adaptadas al tama�o del grupo y su dispersi�n a trav�s de la red. Para cada especificaci�n, puede haber m�ltiples implementaciones, cada una de las cuales contiene el mismo \textit{ModuleSpecID}.


\subsection{Mensajes}

Los servicios y aplicaciones JXTA se comunican usando mensajes JXTA, que son las unidades b�sicas de intercambio de datos entre pares. Cada protocolo JXTA est� definido como un conjunto de mensajes que intercambian los pares participantes. Los mensajes son enviados entre los pares utilizando el servicio de destino (\textit{Endpoint Service}), el servicio de tuber�a (\textit{Pipe Service}), as� como un socket y otras v�as de acceso. La mayor�a de las aplicaciones no necesitan utilizar directamente tuber�as unidireccionales o el servicio de destino, sino que utilizan com�nmente \textit{sockets} de comunicaci�n y tuber�as bidimensionales.
 
Los protocolos de JXTA est�n especificados como un conjunto de mensajes intercambiados entre pares. El uso de mensajes de XML para definir los protocolos permite a muchos tipos diferentes de pares utilizar un protocolo determinado. Debido a que los datos est�n etiquetados, cada par es libre para implementar el protocolo de un modo m�s adecuado a sus capacidades y rol. Por ejemplo, un par que est� muy limitado y no tenga capacidad suficiente para procesar la totalidad de un mensaje, puede usar las etiquetas de datos para extraer las partes que pueda procesar e ignorar las restantes. 

\subsection{Tuber�as (\textit{Pipes})}

Los pares de JXTA utilizan tuber�as para enviar mensajes a otro par. Por tanto, las tuber�as son un mecanismo de transferencia de mensajes as�ncrono, unidireccional y no confiable (con la excepci�n de las tuber�as punto a punto segurass) usados para la comunicaci�n y transferencia de datos. Las tuber�as son canales de comunicaci�n virtuales y pueden conectar pares que no tienen un vinculo f�sico directo, lo que se traduce en una conexi�n l�gica. Las tuber�as pueden ser usadas para enviar cualquier tipo de datos incluyendo texto XML y HTML, im�genes, m�sica, c�digo binario, cadenas de datos y objetos Java.

La tuber�as de destino o de ''punto final'' (\textit{endpoints}) est�n din�micamente vinculadas con los pares de destino (correspondientes a interfaces de red, como un puerto TCP) cuando la tuber�a se abre. Toda la resoluci�n y comunicaci�n de tuber�as se realiza dentro del alcance de un grupo de pares. es decir, las tuber�as de salida y entrada deber�n pertenecer al mismo grupo de pares.
 
Las tuber�as ofrecen dos modos de comunicaci�n: punto a punto (\textit{unicast}) y propagaci�n, como se puede observar en la Figura \ref{fig:pipes}.

\imagen{./imagenes/tuberias}{0.7}{Tipos de tuber�as}{fig:pipes}
 
\begin{milista}
	\item \textbf{Tuber�as punto a punto}: una tuber�a punto a punto conecta exactamente dos tuber�as de punto final. La tuber�a de entrada en un par recibe mensajes enviados desde la tuber�a de salida del otro par. Tambi�n es posible para varios pares unirse a una sola tuber�a de entrada.
 	\item \textbf{Tuber�as de propagaci�n}: este tipo de tuber�as conecta una tuber�a de salida a varias tuber�as de entrada. Los mensajes fluyen desde la tuber�a de salida, que es la fuente de propagaci�n, hacia las tuber�as de entrada.
 	\item \textbf{Tuber�as punto a punto seguras}: es un tipo de tuber�a punto a punto que provee un canal de comunicaci�n seguro y fiable.
\end{milista}
 

\subsection{Anuncios (\textit{Advertisements})}

Todos los recursos de la red JXTA (servicios, pares, grupos de pares y tuber�as) son representados como anuncios. Los anuncios son un lenguaje de estructuras de meta-datos representados como documentos XML (ver Figura \ref{fig:anuncio}). Los protocolos JXTA usan anuncios para describir y publicar la existencia de recursos de un par. Del mismo modo, los pares descubren recursos mediante la b�squeda de sus correspondientes anuncios y pueden guardar en cach� cualquier anuncio local descubierto.
 
Cada anuncio es publicado con un tiempo de vida que especifica la disponibilidad de sus recursos asociados. El tiempo de vida permite la supresi�n de los recursos obsoletos sin la necesidad de un control centralizado, de modo que un anuncio puede ser publicado de nuevo (antes de que el anuncio original expire) para extender la vida �til de un recurso.
 
Los protocolos JXTA definen los siguientes tipos de anuncios:

\begin{milista}
	\item \textbf{Anuncios de pares (\textit{Peer Group Advertisement})}: describen los recursos de un par. El principal uso de estos anuncios es para mantener informaci�n espec�fica sobre el par, como su nombre, su identificador, destinos disponibles, y cualquier atributo que se desee publicar en tiempo de ejecuci�n.
	\item \textbf{Anuncios de grupos de pares (\textit{Peer Group Advertisement})}: describe los recursos de un grupo espec�fico de pares, como son el nombre, el identificador del grupo, descripci�n, etc.
	\item \textbf{Anuncios de tuber�a (\textit{Pipe Advertisement})}: describe un canal de comunicaci�n de una tuber�a, y es usado por el servicio de tuber�a para crear las tuber�as asociadas a una entrada y a una salida. Cada anuncio de tuber�a contiene un identificador simb�lico y un tipo de tuber�a (punto a punto, propagada, segura). 
		\item \textbf{Anuncio de clases de un m�dulo (\textit{Module Class Advertisement})}: describe una clase de un m�dulo. Su objetivo principal es documentar formalmente la existencia de una clase del m�dulo. Incluye un nombre, descripci�n, y un �nico identificador (\textit{ModuleClassID}).
		\item \textbf{Anuncio de especificaci�n de un m�dulo (\textit{Module Specification Advertisement})}: describe una especificaci�n de un m�dulo. Su principal objetivo es proveer referencias a la documentaci�n necesaria con el fin de crear implementaciones conforme con la especificaci�n. Un uso secundario es hacer correr instancias que se puedan utilizar de manera remota mediante la publicaci�n de informaci�n. Este anuncio incluye nombre, descripci�n, identificador �nico de la especificaci�n del m�dulo (\textit{ModuleSpecID}), anuncio de tuber�a y un campo que contiene par�metros arbitrarios a ser interpretados por cada aplicaci�n.
	\item \textbf{Anuncios de implementaci�n de un m�dulo (\textit{Module Implementation Advertisement})}: define una implementaci�n de una especificaci�n de un m�dulo. Incluye un nombre, el \textit{ModuleSpecID} asociado, as� como el c�digo, paquete y campos de par�metros que permiten a un par recuperar los datos necesarios para ejecutar la implementaci�n.
	\item \textbf{Anuncio de citas (\textit{Rendezvous Advertisement})}: describe como se establecen citas entre pares de un grupo.
	\item \textbf{Anuncio de informaci�n de un par (\textit{Peer Info Advertisement})}: describe los recursos de informaci�n de un par. El uso principal de este anuncio es contener informaci�n especifica acerca del estado actual de un par, tales como el tiempo de actividad, cuenta de mensajes de entrada y salida, el tiempo que ha pasado desde el ultimo mensaje recibido y el tiempo que ha pasado desde el ultimo mensaje enviado.
\end{milista}
  
\imagen{./imagenes/anuncio}{0.6}{Ejemplo de un anuncio de una tuber�a}{fig:anuncio}
 

\subsection{Seguridad}

Redes P2P din�micas, como la red JXTA, necesitan soportar diferentes niveles de acceso a los recursos. En JXTA, los pares operan en un modelo basado en roles, en el cual un par individual act�a bajo la autoridad de otro par de confianza para realizar una tarea.

Se proveen cinco requisitos b�sicos de seguridad en JXTA: 

\begin{milista}
	\item \textbf{Confidencialidad}: garantiza que el contenido de un mensaje no se divulgar� a pares no autorizados.
	\item \textbf{Autenticaci�n}: garantiza que el emisor de un mensaje es qui�n dice ser.
	\item \textbf{Autorizaci�n}: garantiza que el emisor est� autorizado para enviar un mensaje.
	\item \textbf{Integridad de los datos}: garantiza que el mensaje no ha sido modificado accidental o deliberadamente mientras se transmit�a.
	\item \textbf{Refutabilidad}: garantiza que el mensaje fue transmitido por un emisor identificado y que no es una respuesta de un mensaje previamente transmitido.
\end{milista}

Los mensajes XML permiten a�adir nuevos meta-datos como credenciales, certificados y claves p�blicas a los mensajes de JXTA, permitiendo que se cumplan los requisitos de seguridad anteriores. Adem�s, los mensaje se pueden encriptar y firmar.

Un credencial es un token que se utiliza para identificar un emisor y para comprobar que tiene autorizaci�n para enviar el mensaje. Adem�s, este token est� presenta en cada mensaje que ese emisor env�a, comprob�ndose con la identidad del emisor cuando se recibe el mensaje. Cabe destacar que la implementaci�n de las credenciales se realiza como una configuraci�n de un plug-in, por lo que pueden exisitr m�ltiples configuraciones de credenciales dentro de la misma red. Con esto se intenta que JXTA sea compatible con otros mecanismos de seguridad como SSL (\textit{Secure Socket Layer)} e IPSec (\textit{Internet Protocol Secure}).
\clearpage
% Estos puntos se sacan de la gu�a de programaci�n y de la secci�n "JXTA Overview" del Nutshell
\section{Arquitectura de JXTA}

\subsection{Componentes de JXTA}
% Incluir la secci�n "aspectos clave"

\subsection{Arquitectura en red de JXTA}


\clearpage
\section{Protocolos utilizados en JXTA}

\subsection{Protocolo de Descubrimiento de Pares (\textit{PDP - Peer Discovery Protocol})}

\subsection{Protocolo de Informaci�n de Pares (\textit{PIP - Peer Information Protocol})}

\subsection{Protocolo de Resoluci�n de Pares (\textit{PRP - Peer Resolver Protocol})}

\subsection{Protocolo de Enlace a Tuber�as (\textit{PBP - Pipe Binding Protocol})}

\subsection{Protocolo de Encaminamiento al Destino (\textit{ERP - Endpoint Routing Protocol})}

\subsection{Protocolo de Citas (\textit{RVP - Rendezvous Protocol})}
\clearpage
\section{Iniciaci�n a la programaci�n de aplicaciones con JXTA}

El objetivo de esta secci�n es que sirva al lector como gu�a de referencia en su iniciaci�n a la programaci�n de aplicaciones basadas en JXTA. Para ello se comenzar� configurando el entorno de Java y se dar�n unas nociones b�sicas de configuraci�n del entorno de JXTA. Acto seguido, se indicar� como ejecutar el \textit{shell} de JXTA y, por �ltimo, se har� un peque�o ejemplo tipo \textit{HelloWorld}.

\subsection{Requisitos del sistema}

La implementaci�n actual de JXTA requiere que el sistema operativo tenga instalado una versi�n igual o posterior a la 5.0 del entorno de ejecuci�n de Java (JRE o \textit{Run-Time Environtment}) o del entorno de desarrollo (SDK o \textit{Software Development Kit}). Tanto el JDK como el JRE est�n disponibles para sistemas basados en Solaris, Microsoft Windows, Linux, Mac OS X, etc�tera.

\subsection{Instalaci�n y configuraci�n de los entornos}
\subsubsection{Java}
Para comprobar si ya dispone de una instalaci�n del entorno de ejecuci�n de Java en el equipo y que est� configurado correctamente, abra una ventana de comandos \textit{(Inicio -\textgreater Ejecutar -\textgreater cmd)} y escriba "`java -version"'. Debe obtener un resultado similar al de la figura \ref{fig:java-version}:

\imagen{./imagenes/java-version.png}{0.80}{C�mo consultar la versi�n de Java instalada}{fig:java-version}

Si aparece un mensaje de error, pruebe lo siguiente: 
\begin{milista}
	\item Compruebe si tiene instalado Java. Para ello compruebe que exista el directorio Java en \textit{C:\textbackslash Archivos de programa} (t�picamente). Si lo est�, config�relo como se indica m�s abajo.
	\item Si no existe, vaya a la siguiente direcci�n de Internet y descargue la �ltima versi�n del  Java SE Development Kit (JDK) \cite{download-jdk}
\end{milista}

Una vez descargado, ejecute el archivo, elija el directorio de instalaci�n y siga los pasos para completar la instalaci�n. Tras instalarlo, debe configurarlo. Para ello vaya a \textit{Inicio} y haga clic derecho en \textit{Mi PC}. En las opciones, elija la pesta�a de \textit{Opciones Avanzadas} y haga clic en \textit{Variables de entorno}. En la nueva ventana que aparece, busque la variable \textit{Path} y a�ada al final \texttt{;C:\textbackslash Archivos de programa\textbackslash Java\textbackslash jdk1.x.x\_xx\textbackslash bin}, donde "`x"' hace referencia a la versi�n que tiene instalada de Java. Ese es el directorio de instalaci�n por defecto de Java. Si en la instalaci�n eligi� otro diferente, debe introducir el que eligi� al instalarlo.

Una vez realizados los pasos descritos, pruebe a ejecutar de nuevo la sentencia "`java -versi�n"' en una ventana de comandos.


\subsubsection{JXTA}
Para comenzar, debe hacerse de una versi�n funcional de la librer�a de JXTA. Se proponen dos m�todos (ambos requieren que visite la p�gina de descargas de JXTA \cite{download-jxta}):
\begin{milista}
	\item Descargue el fichero \textit{jxse-src-2.5.zip}. Descompr�malo y comp�lelo con \textit{ant} o \textit{make}. Acto seguido se generar� el fichero \textit{jxta.jar} en el directorio \textit{./jxse-src-2.5/dist/}.
	\item Descargue el fichero \textit{jxse-lib-2.5.zip}. En su interior se encuentra el fichero \textit{jxta.jar} ya compilado.
\end{milista}

Tambi�n puede descargar la documentaci�n de JXTA (fichero \textit{jxse-doc-2.5.zip}) o la gu�a de programaci�n de JXTA para Java SE/EE (JXSE) 2.5 \cite{jxta-programming-guide}.

En versiones anteriores era necesario realizar una instalaci�n de dicho motor, pero hoy en d�a basta con incluir la librer�a de JXTA en la aplicaci�n y realizar una configuraci�n del \textit{peer}, bien sea en el c�digo fuente de la aplicaci�n o mediante un \textit{configurator} externo.



%JXSE includes simple UI configurator. In most cases this default may be overridden by using a programmatic platform configurator such as the NetworkConfigurator class, or by using the NetworkManager, which abstracts configuration to one of the preset configurations:
%1.Ad-Hoc: A node which is typically deployed in an ad-hoc network. The peer will not use any infrastructure peers (Rendezvous or Relay services).
%2.Edge: In addition to supporting the Ad-Hoc behavior, an Edge node can attach to an infrastructure peer (a Rendezvous, Relay, or both).
%3.Rendezvous: Provides network bootstrapping services, such as discovery, pipe resolution, etc.
%4.Relay: Provides message relaying services, enabling cross firewall traversal.
%5.Proxy: Provides JXME JXTA for J2ME proxy services.
%6.Super: Provides the functionality of a Rendezvous, Relay, and Proxy node.
%It is highly recommended that JXTA peers be configured through the NetworkManager, or NetworkConfigurator static configuration methods with the desired mode of operation such as newEdgeConfiguration(storeHome).
%Once a configuration is created, it maybe tuned to the application requirements.

% Secci�n de ap�ndices
% Se a�ade el titulo de "Ap�ndices" al �ndice
%\addcontentsline{toc}{section}{Ap�ndices}

%\input{Apendices.tex}

% A�adimos la bibliografia al �ndice
\clearpage\phantomsection
\addcontentsline{toc}{section}{Referencias}

% Bibliograf�a. En este caso se usa BibTeX
\bibliographystyle{plain}
\pagestyle{plain} 
\bibliography{Bibliografia}

\end{document}
